\subsection{同じPI次数を保つ中央商の中心持上げ}
\label{subsec:cubic-third-root-unit-boundary-L-trace-center}
\begin{lemma}[trace-closed中央商]
\label{lem:cubic-trace-closed-central-reduction}
$T$ は標数0体上のprime PI代数で、PI次数を $n$ とする。
$T$ は自身のtrace ringに等しいと仮定する。
$x\in Z(T)$ は正則で、$S=T/(x)$ はnoetherian整域かつPI次数 $n$ とする。
このとき自然な写像 $Z(T)\to Z(S)$ は全射で、核は $xZ(T)$ である。
\end{lemma}
\begin{proof}
$R=\operatorname{im}(Z(T)\to Z(S))$、$L=\operatorname{Frac}R$ と置く。
trace-closed性から、各 $a\in T$ の次数 $n$ のreduced characteristic polynomialは
$Z(T)$ に係数を持つ。従って任意の $s\in S$ は $R$ 上次数 $n$ のmonic多項式を満たす。
$S_L$ の全元も $L$ 上次数 $n$ 以下で代数的である。
非零元の消滅多項式から $t$ の因子を除けば定数項を非零にできるので、
その逆元は同じ元の多項式で与えられる。
従って $S_L$ は除法環であり、$Q(S)$ に等しい。

$K=Z(Q(S))$ の各元の $L$ 上の次数も $n$ 以下である。
標数0のprimitive element定理から、$[K:L]$ は有限で $n$ 以下。
実際、各有限部分拡大は単拡大となり次数 $n$ 以下なので、全拡大の次元も $n$ 以下となる。
$Q(S)$ の $K$ 上の次数は仮定により $n$。
その最大部分体 $M/K$ は次数 $n$ であり、$M/L$ は有限分離拡大である。
再びprimitive elementを取ると、その次数は $n[K:L]$ であるが、
$Q(S)=S_L$ の元の次数上界 $n$ により $[K:L]=1$。
従って $K=L$ である。

$T$ のreduced traceは $xT$ を $xZ(T)$ に写すので、
$R$ 線型な巡回写像 $\bar t:S\to R$ を誘導し、$\bar t(1)=n$ を満たす。
これを $L$ へ局所化すると、中心単純代数 $Q(S)$ 上の $L$ 線型巡回写像になる。
このような写像はreduced traceのscalar倍であり、単位元の値からそれ自身に等しい。
従って任意の $z\in Z(S)$ について $\bar t(z)=nz\in R$、すなわち $z\in R$ となる。
最後に、$xa\in Z(T)$ なら $x[a,b]=0$ が全 $b\in T$ で成り立ち、
$x$ の正則性から $a\in Z(T)$。これが核の記述を与える。
\end{proof}

\begin{theorem}[全 $L/\Omega$ 上の中心持上げ]
\label{thm:cubic-third-root-unit-boundary-L-trace-center}
本族の全楕円パラメータと全 $u,v$ に対し
\[
 Z(C)/(x)\simeq Z(S_C),\qquad
 \operatorname{gr} Z(D)=Z(S_C).
\]
$C$ は全中心上有限で、その中心は正規である。
ここでは各幾何学的点についての主張であり、任意底変更での全中心形成までは主張しない。
\end{theorem}
\begin{proof}
既証の正規一次拡大認識とLBS\cite[Corollaries~2.3, 2.7]{LBS}から、
$C$ はnoetherian、Auslander-regular、Cohen--Macaulayである。
またconnected gradedかつ有限大域次元なのでstably freeである。
後者は、有限生成projectiveにgood filtrationを入れ、
associated gradedの有限graded free resolutionをfiltered free resolutionへ持ち上げ、
その分裂を交互に取ることで直接分かる。
前節でPI次数9を証明した $D$ と $C$ は、中央Laurent局所化で同じPI次数を持つ。
従ってStafford\cite[Theorem~2.10, Corollary~2.11]{StaffordMaxOrders}により、
$C$ は自身のtrace ringに等しく、全中心上有限、中心は正規となる。
$S_C=C/(x)$ も整域でPI次数9なので、前補題を適用できる。
中心は次数付きであるから、斉次中心元は同じ次数の中心元へ持ち上がる。
中央 $x$ の局所化の次数零部分を取れば $D$ の濾過中心の主張が従う。
\end{proof}
